% !Mode:: "TeX:UTF-8"
% 第 1 章，照《研究生学位论文书写范例（理工类）》第 1 章节选，省略处写“……”。
% \chapter、\section 等方括号里是英文题目：博士学位论文的英文目录要它，硕士、本科
% 写了也不排，所以一律写上最省事。

\chapter{绪论}[Introduction]

\section{课题背景及研究的目的和意义}[Background, objective and significance of the subject]

发展国防工业、微电子工业等尖端技术需要精密和超精密的仪器设备，精密仪器设备要求高速、……

……

\section{气体润滑轴承及其相关理论的发展概况}[Developmental of gas-lubricated bearing and correlated theories]

气体轴承是利用气膜支撑负荷或减少摩擦的机械构件。……

……

\subsection{气体润滑轴承的发展}[Developmental of gas-lubricated bearing]

1828年，R.R.Willis\cite{zhao1998}发表了一篇关于小孔节流平板中压力分布的文章，这是有记载的研究气体润滑的最早文献。……

根据间隙内气膜压力的产生原理，气体轴承可以分为四种基本形式：

（1）气体静压轴承\quad 加压气体经过节流器进入间隙，在间隙内产生压力气膜使物体浮起的气体轴承，……

\subsection{气体润滑轴承的分类}[Classification of gas-lubricated bearing]

根据间隙内气膜压力的产生原理，气体轴承可以分为四种基本形式，其结构如图\ref{fig:1-1}所示。

% \hitindex{zh=…, en=…} 排出当前语种的那个词，并往中英文两个索引各登记一条。
（1）\hitindex{zh=气体静压轴承, en=Externally pressurized gas bearing}\quad 加压气体经过节流器进入间隙，在间隙内产生压力气膜使物体浮起的气体轴承，结构如图\ref{fig:1-1}(a)所示。……

……

% 四张分图、分图题写在图题下面，照范例。博士学位论文的图题中英文两行，英文写在方括号里；
% 硕士、本科只排中文，方括号照写也不碍事。
\begin{figure}[!ht]
  \centering
  \subfigure[]{\includegraphics[width=6.92cm]{aerostatic}}%
  \hspace{1cm}%
  \subfigure[]{\includegraphics[width=6.65cm]{aerodynamic}}

  \subfigure[]{\includegraphics[width=6.72cm]{hybrid}}%
  \hspace{0.5cm}%
  \subfigure[]{\includegraphics[width=7.71cm]{squeeze-film}}
  \caption{气体润滑轴承的分类}[Classification of the gas-lubricated bearing]
  \label{fig:1-1}

  {\wuhao (a) 气体静压轴承；(b) 气体动压轴承；(c) 气体动静压轴承；(d) 气体压膜轴承\par}
\end{figure}

（也可以按照下图范例书写，分图题写在各自的分图下面）

% 一张图分两页排时，后一页那半张写成续图：\ContinuedFloat 让分图接着编 (c)(d)。
\begin{figure}[!ht]
  \centering
  \subfigure[气体静压轴承]{\label{fig:1-2a}\includegraphics[width=6.87cm]{aerostatic}}%
  \hspace{1cm}%
  \subfigure[气体动压轴承]{\label{fig:1-2b}\includegraphics[width=6.03cm]{aerodynamic}}
  \caption{气体润滑轴承的分类}[Classification of the gas-lubricated bearing]
  \label{fig:1-2}
\end{figure}

\clearpage

\begin{figure}[!ht]
  \ContinuedFloat
  \centering
  \subfigure[气体动静压轴承]{\includegraphics[width=6.70cm]{hybrid}}%
  \hspace{1cm}%
  \subfigure[气体压膜轴承]{\includegraphics[width=7.17cm]{squeeze-film}}
  \caption*{图~\thefigure（续图）}
\end{figure}

\subsection{多孔质气体静压轴承的研究}[Research on porous externally pressurized gas bearing]

由于气体的压力低和可压缩性，……。

\subsubsection{多孔质静压轴承的分类}[Classification of porous thrust bearing]

轴承工作面的整体或……。

\subsubsection{多孔质材料特性的研究}[Research on characteristics of porous materials]

材料的主要特点是具有一定的……。

（1）孔隙特性\quad 多孔质材料是由……。

……

\section{本文的主要研究内容}[Main research content of this subject]

本课题的研究内容主要是针对局部多孔质止推轴承的多孔质材料的渗透率、静压轴承的静态特性、稳定性及其影响因素进行展开，……。
